% 作者题证的编辑转录副本，不是独立下载的上游原件。
% 来源：https://github.com/sanathdevalapurkar/algtop-notes/blob/3f5d3189e2082716a69fccc1711d02ed848552d2/906/lec-68-leray-hirsch.tex

\paragraph{Filtration used in the proof.}
Let $\pi:E\to B$ be a fibration, with path-connected base and fibre $F$.
Cohomology has coefficients in a commutative ring $R$. The multiplicative Serre
spectral sequence is
\[
 E_2^{s,t}=H^s(B;\underline{H^t(F;R)})\Longrightarrow H^{s+t}(E;R).
\]
The pullback $\pi^*$ makes $H^*(E;R)$ an $H^*(B;R)$-module. On each total
degree, rewrite the decreasing Serre filtration as
$F_tH^n(E;R)=F^{n-t}H^n(E;R)$. Then
\[
 F_0H^n(E;R)=E_\infty^{n,0}=\im(\pi^*:H^n(B;R)\to H^n(E;R)),
 \qquad \operatorname{gr}_tH^*(E;R)=E_\infty^{*,t}.
\]
This is an exhaustive, bounded-below filtration by $H^*(B;R)$-modules, and
the displayed identification is an identification of modules.

\begin{theorem}[Leray--Hirsch]
Suppose $B$ and $F$ are path connected, each $H^t(F;R)$ is a finite-rank free
$R$-module, and restriction $i^*:H^*(E;R)\to H^*(F;R)$ is surjective.
Choose a degree-preserving $R$-linear section
$\sigma:H^*(F;R)\to H^*(E;R)$ of $i^*$. Then
\[
 \overline\sigma:H^*(B;R)\otimes_RH^*(F;R)\longrightarrow H^*(E;R),
 \qquad x\otimes y\longmapsto\pi^*(x)\smile\sigma(y)
\]
is an isomorphism of $H^*(B;R)$-modules.
\end{theorem}
\begin{remark}
The isomorphism depends on the choice of $\sigma$. In particular, the theorem
asserts freeness as an $H^*(B;R)$-module, not a canonical algebra isomorphism.
\end{remark}
\begin{proof}
Use the Serre spectral sequence above. Restriction to the fibre is an edge
homomorphism and factors through
\[
 H^*(E;R)\longrightarrow E_\infty^{0,*}
 \hookrightarrow E_2^{0,*}=H^0(B;\underline{H^*(F;R)})
 \hookrightarrow H^*(F;R).
\]
Surjectivity of restriction implies
$H^0(B;\underline{H^*(F;R)})=H^*(F;R)$. Thus the monodromy action of
$\pi_1(B)$ on $H^*(F;R)$ is trivial. Since each graded piece of fibre
cohomology is finite free, the second page is
\[
 E_2^{s,t}=H^s(B;R)\otimes_R H^t(F;R).
\]
There are no outgoing differentials on the fibre column: otherwise the
restriction to the fibre could not be surjective. There are no outgoing
differentials on the base row either. The differential is a derivation,
and this tensor-product algebra is generated by the base row and fibre
column. Consequently all differentials vanish, page by page. Hence
\[
 E_
\infty^{*,t}=H^*(B;R)\otimes_RH^t(F;R).
\]
Filter the tensor product by degree in fibre cohomology:
\[
 F_q\bigl(H^*(B;R)\otimes_RH^*(F;R)\bigr)
   =\bigoplus_{t\leq q}H^*(B;R)\otimes_RH^t(F;R).
\]
The map $\overline\sigma$ preserves this filtration and the corresponding
filtration on $H^*(E;R)$. On the associated graded its map is the
identification
\[
 H^*(B;R)\otimes_RH^t(F;R)
   =E_\infty^{*,t}=\operatorname{gr}_tH^*(E;R).
\]
It is therefore an isomorphism on the associated graded. Since the
filtrations are exhaustive and bounded below, $\overline\sigma$ itself is
an isomorphism.
\end{proof}
